% -------------------------------------------------------------- %
%      csd5748 - HY-225 - Spring 2025 - CSD-UOC - Ex. 13         %
% -------------------------------------------------------------- %

\documentclass{article}

\usepackage[utf8]{inputenc} 
\usepackage[english, greek]{babel}
\usepackage{subcaption}
\usepackage{hyperref}
\usepackage{xcolor}
\hypersetup{ hidelinks, colorlinks=true, linkcolor=blue, urlcolor=blue }
\captionsetup[subfigure]{labelformat=empty}
\usepackage{amsmath}
\usepackage{amsfonts}
\usepackage{amssymb}
\usepackage{array}
\usepackage{graphicx}
\usepackage[a4paper, left=1.1cm, right=1.1cm, top=2cm, bottom=2cm]{geometry}
\usepackage[absolute, overlay]{textpos}
\usepackage{placeins}
\usepackage{listings}
\usepackage{xcolor}
\setlength{\jot}{0pt}
\setlength{\mathsurround}{0.1cm}
\setlength{\jot}{0.1cm}
\newcommand{\skato}{\hspace{0.25cm}}
\newcommand{\skata}{\hspace{1cm}}
\usepackage[english, greek]{babel}  
\usepackage{tcolorbox} 
\usepackage{titlesec}
\usepackage{fancyvrb} 
\geometry{margin=2.5cm}
\lstset{
  language=C,
  inputencoding=utf8,
  extendedchars=true,
  basicstyle=\ttfamily\small,
  keywordstyle=\color{blue},
  commentstyle=\color{gray}\itshape,
  stringstyle=\color{red},
  showstringspaces=false,
  columns=fullflexible,
  keepspaces=true,
  frame=single,
  breaklines=true,
  tabsize=2,
  escapeinside={(*@}{@*)}, 
  morekeywords={char, return, if, while, volatile, unsigned, int, define},
}

\begin{document}

% -------------------------------------------------------------- %

\text{   }
\vspace{-1.0cm}
\begin{center}
\setlength{\fboxrule}{2pt}
\setlength{\fboxsep}{15pt} 
\fbox
{
{\noindent{\Huge \textbf{\textlatin{HY225 -} \textbf{Ασκήσεις 13} \textlatin{- csd5748}}}}
}
\end{center}
\vspace{3cm}


% -------------------------------------------------------------- %


\begin{center}
    {\underline{\Large \textbf{Άσκηση 13.4}}}
\end{center}
\vspace{0.12345cm}


\section*{\text{α})}

\setlength{\fboxrule}{1pt}
\begin{figure}[h]
  \centering
  \fbox{\includegraphics[width=0.75\textwidth]{13.jpg}}
  \caption*{$*$ Θεωρητικό Κύκλωμα Σχεδιασμένο στο \href{https://circuitverse.org/simulator}{\textlatin{CircuitVerse}}.}
\end{figure}



\vspace{0.1cm}
\section*{\text{β})}

\[
\begin{tabular}{|c|c|c|}
\hline
\textlatin{Hex Address} & \textlatin{Binary Address} & \text{Θέση Μνήμης \textlatin{(Decimal)}} \\
\hline
\textlatin{000} & \textlatin{00000000000} & \textlatin{Main Memory \ \ \ \ \ 0} \\
\textlatin{00E} & \textlatin{00000001110} & \textlatin{Main Memory \ \ \ 14} \\
\textlatin{100} & \textlatin{00100000000} & \textlatin{Main Memory \ \ 256} \\
\textlatin{200} & \textlatin{01000000000} & \textlatin{Main Memory \ \ 512} \\
\textlatin{2FF} & \textlatin{01011111111} & \textlatin{Main Memory \ \ 767} \\
\textlatin{300} & \textlatin{01100000000} & \textlatin{Main Memory \ \ 768} \\
\textlatin{3FF} & \textlatin{01111111111} & \textlatin{Main Memory \ 1023} \\
\textlatin{400} & \textlatin{10000000000} & \textlatin{Large I/O \ \ \ \ 0} \\
\textlatin{4FF} & \textlatin{10011111111} & \textlatin{Large I/O \ 255} \\
\textlatin{500} & \textlatin{10100000000} & \textlatin{Large I/O \ 256} \\
\textlatin{5FF} & \textlatin{10111111111} & \textlatin{Large I/O \ 511} \\
\textlatin{600} & \textlatin{11000000000} & \text{Παράνομη} \\
\textlatin{608} & \textlatin{11000001000} & \text{Παράνομη} \\
\textlatin{60A} & \textlatin{11000001010} & \text{Παράνομη} \\
\textlatin{60E} & \textlatin{11000001110} & \text{Παράνομη} \\
\textlatin{60F} & \textlatin{11000001111} & \text{Παράνομη} \\
\textlatin{610} & \textlatin{11000010000} & \text{Παράνομη} \\
\textlatin{680} & \textlatin{11010000000} & \text{Παράνομη} \\
\textlatin{6C0} & \textlatin{11011000000} & \text{Παράνομη} \\
\textlatin{6C3} & \textlatin{11011000011} & \text{Παράνομη} \\
\textlatin{6CF} & \textlatin{11011001111} & \text{Παράνομη} \\
\textlatin{6FF} & \textlatin{11011111111} & \text{Παράνομη} \\
\textlatin{700} & \textlatin{11100000000} & \text{Παράνομη} \\
\textlatin{708} & \textlatin{11100001000} & \text{Παράνομη} \\
\textlatin{70A} & \textlatin{11100001010} & \text{Παράνομη} \\
\textlatin{70E} & \textlatin{11100001110} & \text{Παράνομη} \\
\textlatin{70F} & \textlatin{11100001111} & \text{Παράνομη} \\
\textlatin{710} & \textlatin{11100010000} & \text{Παράνομη} \\
\textlatin{780} & \textlatin{11110000000} & \text{Παράνομη} \\
\textlatin{7C0} & \textlatin{11111000000} & \text{\textlatin{IN Status} \ \ \ 1984} \\
\textlatin{7C3} & \textlatin{11111000011} & \text{\textlatin{\ OUT Data} \ 1987} \\
\textlatin{7CF} & \textlatin{11111001111} & \text{Παράνομη} \\
\textlatin{7FF} & \textlatin{11111111111} & \text{Παράνομη} \\
\textlatin{800} & \textlatin{100000000000} & \text{Παράνομη (ανύπαρκτη)} \\
\textlatin{808} & \textlatin{100000001000} & \text{Παράνομη (ανύπαρκτη)} \\
\textlatin{80A} & \textlatin{100000001010} & \text{Παράνομη (ανύπαρκτη)} \\
\textlatin{80E} & \textlatin{100000001110} & \text{Παράνομη (ανύπαρκτη)} \\
\textlatin{80F} & \textlatin{100000001111} & \text{Παράνομη (ανύπαρκτη)} \\
\textlatin{C0E} & \textlatin{110000001110} & \text{Παράνομη (ανύπαρκτη)} \\
\textlatin{C0F} & \textlatin{110000001111} & \text{Παράνομη (ανύπαρκτη)} \\
\textlatin{F02} & \textlatin{111100000010} & \text{Παράνομη (ανύπαρκτη)} \\
\textlatin{F12} & \textlatin{111100010010} & \text{Παράνομη (ανύπαρκτη)} \\
\hline
\end{tabular}
\]




% ---------------------------------------------------------------------------------------



\newpage

\begin{center}
    {\underline{\Large \textbf{Άσκηση 13.5}}}
\end{center}

\bigskip
\bigskip

\selectlanguage{english}

\begin{lstlisting}
char read_kbd_busywait_char()
{
    volatile unsigned long int *status = (volatile unsigned long int *) 0x7C0;
    volatile unsigned long int *data = (volatile unsigned long int *) 0x7C3;
    while (*status != 1) {};
    return ((char) *data);
}
\end{lstlisting}

\bigskip

\begin{lstlisting}
char read_kbd_polling_char()
{
    volatile unsigned long int *status = (volatile unsigned long int *) 0x7C0;
    volatile unsigned long int *data = (volatile unsigned long int *) 0x7C3;
    if (*status_reg == 1) { return ((char) *data); }
    return ("\0");
}
\end{lstlisting}

\selectlanguage{greek}







\newpage

\text{   }
\begin{center}
    \vspace{-1.5cm}
    {\underline{\Large \textbf{Άσκηση 13.6}}}
\end{center}
\vspace{0.345cm}
 

\subsection*{\text{Δεδομένα (όλα):}}

\begin{itemize}
    \item Συχνότητα \textlatin{CPU}: $1$~\textlatin{GHz} $=$ 1 κύκλοι/\textlatin{ns} $=$ 1.000.000.000 κύκλοι/δευτερόλεπτο.
    \item \textlatin{Overhead} (κόστος) ανά διακοπή: $2000$ κύκλοι.
    \item Κόστος δειγματοληψίας μίας συσκευής: $200$ κύκλοι.
    \item Αριθμός συσκευών/σημάτων: $40$ (μέγιστο ρυθμό 1 \textlatin{kHz}).
    \item Συχνότητα \textlatin{polling}: $1$ διακοπή κάθε $1~\textlatin{ms} = 1000$ διακοπές/δευτερόλεπτο.
    \item Μέσος ρυθμός εισόδων: 1 \& 40 \& 400 είσοδοι/δευτερόλεπτο ανά γραμμή (ανάλογα).
    \item Ε/Ε βάσει διακοπών: Κάθε νέα είσοδος προκαλεί 1 διακοπή.
    \item Ταχύτητα εισόδου: $32~\textlatin{Mbit/sec} = 4~\textlatin{MBytes/sec}$ \& 1~\textlatin{Gbit/second} = 125 \textlatin{MBytes/s.}
    \item Κάθε πακέτο έχει μέγεθος: $40~\textlatin{Bytes}$ \textlatin{(IP packet).}
    \item Ρυθμός άφιξης πακέτων: ${4~(\textlatin{MBytes/sec})}/{40~(\textlatin{Bytes})} = 100.000~\text{πακέτα/δευτερόλεπτο}$.
\end{itemize}

\vspace{0.234567cm}

\subsection*{\text{(α\textlatin{i})}}
\begin{itemize}
\item \text{Κόστος απλά διακοπής} = \textlatin{Overhead} = 2000 \text{ κύκλοι}.
\end{itemize}

\subsection*{\text{(α\textlatin{ii})}}
\begin{itemize}
\item  \text{Κόστος διακοπής} = \textlatin{Overhead} + 40\text{ συσκ.} $\times$ 200\text{ κύκλοι/συσκ.} = 2000 + 8000 = 10.000\text{ κύκλοι}.
\end{itemize}

\subsection*{\text{(α\textlatin{iii})}} 
\begin{itemize}
\item \text{Κόστος ανά δευτ.} = 1000\text{ διακοπές/δευτερόλεπτο} $\times$ 10{.}000\text{ κύκλοι/διακοπή} = 10{.}000{.}000\text{ κύκλοι/δευτ.}
\end{itemize}

\subsection*{\text{(α\textlatin{iv})}} 
\begin{itemize}
\item \text{Ποσοστό συνολικής υπολογιστικής δυναμικότητας επεξ.} = {10.000.000} / {1.000.000.000} $\times$ 100 = {1\%}
\end{itemize}
 

\subsection*{\text{(β\textlatin{i})}}

\begin{itemize}
    \item \text{Συνολικές νέες είσοδοι ανά δευτερόλεπτο: } $40~\text{γραμμές} \times 40~\text{εισόδους} = 1.600~\text{εισόδοι/δευτερόλεπτο}$.
\end{itemize}

\subsection*{\text{(β\textlatin{ii})}}

\begin{itemize}
    \item \text{Αριθμός διακοπών ανά δευτερόλεπτο: } $1600\text{ εισόδοι} \times 1\text{ διακοπή} = 1.600$ διακοπές/δευτ. περίπου.
\end{itemize}

\subsection*{\text{(β\textlatin{iii})}}

\begin{itemize}
    \item \text{Κόστος ανά δευτερόλεπτο: } $1600~\text{διακοπές} \times 2000~\text{κύκλοι/διακοπή} = 3.200.000~\text{κύκλοι/δευτερόλεπτο}$.
\end{itemize}

\subsection*{\text{(β\textlatin{iv})}}

\begin{itemize}
    \item \text{Ποσοστό συνολικής υπολογιστικής ισχύος επεξ.: } ${3.200.000} / {1.000.000.000} \times 100 = 0.32\%$
\end{itemize}

\subsection*{\text{(β\textlatin{v})}}

\begin{itemize}
    \item \textlatin{Polling} (α) = {1\%} $>$
    \textlatin{Interrupts} (β) = {0.32\%} $\implies$ \text{οπότε συμφέρουν οι διακοπές (β)} ...
\end{itemize}

\newpage
\text{   }
\vspace{-2.205cm}

\subsection*{\text{(γ\textlatin{i})}}

\begin{itemize}
    \item \text{Συνολικές νέες είσοδοι ανά δευτερόλεπτο: } $40~\text{γραμμές} \times 400~\text{εισόδους} = 16.000~\text{σήματα/δευτ.}$
\end{itemize}
 
\subsection*{\text{(γ\textlatin{ii})}}

\begin{itemize}
        \item \text{Αριθμός διακοπών ανά δευτερόλεπτο: } $16000\text{ εισόδοι} \times 1\text{ διακοπή} = 16.000$ διακοπές/δευτ. περίπου.
\end{itemize}

\subsection*{\text{(γ\textlatin{iii})}}

\begin{itemize}
    \item \text{Κόστος ανά δευτερόλεπτο: } $16.000~\text{διακοπές} \times 2000~\text{κύκλοι/διακοπή} = 32.000.000~\text{κύκλοι/δευτ.}$
\end{itemize}

\subsection*{\text{(γ\textlatin{iv})}}

\begin{itemize}
    \item \text{Ποσοστό συνολικής υπολογιστικής ισχύος: } ${32.000.000}/{1.000.000.000} \times 100 = 3.2\%$
\end{itemize}

\subsection*{\text{(γ\textlatin{v})}}

\begin{itemize}
    \item \textlatin{Polling} (α) = {1\%} $>$
    \textlatin{Interrupts} (γ) = {3.2\%} $\implies$ \text{οπότε συμφέρει η δειγματοληψία (α) τώρα} ...
\end{itemize}
 

\subsection*{\text{(δ\textlatin{i})}}

\begin{itemize}
    \item \text{Συνολικό κόστος για 40 διακοπές: } $40 \times 2000\text{ διακοπές} = 80.000$ κύκλοι.
\end{itemize}

\subsection*{\text{(δ\textlatin{ii})}}

\begin{itemize}
    \item \text{Χρόνος για όλους αυτούς: } $80.000\text{ κύκλοι} / 1\text{ (\textlatin{ns}/κύκλο}) = 80.000 \textlatin{ns} = 80 \text{μ\textlatin{s}} \approx 0.00008~\textlatin{secs}$
\end{itemize}

\subsection*{\text{(δ\textlatin{iii})}}

\begin{itemize}
    \item \text{Αφού } 80 \text{μ\textlatin{s}} $<$ 1\textlatin{ms} = 1000 \text{μ\textlatin{s}} $\implies$ \text{προλαβαίνει ο επεξεργαστής και δεν θα χαθεί κανένα σήμα ...}
\end{itemize} 
 

\subsection*{\text{(ε\textlatin{i})}}

\begin{itemize}
    \item \text{Συνολικές νέες είσοδοι (πακέτα) ανά δευτερόλεπτο: } $100.000$ \text{(δεδομένο)}.
\end{itemize}

\subsection*{\text{(ε\textlatin{ii})}}

\begin{itemize}
    \item \text{Διακοπές ανά δευτερόλεπτο: } $100.000 \times 1\text{ διακοπή/πακέτο} = 100.000$.
\end{itemize}

\subsection*{\text{(ε\textlatin{iii})}}

\begin{itemize}
    \item \text{Συνολικοί κύκλοι επεξεργαστή για τις διακοπές: } $100.000 \times 2000\text{ κύκλοι/διακοπή} = 200.000.000$ κύκλοι.
\end{itemize}

\subsection*{\text{(ε\textlatin{iv})}}

\begin{itemize}
    \item \text{Ποσοστό υπολογιστικής ισχύος που καταναλώνεται: } ${200.000.000} / {1.000.000.000} \times 100 = 20\%$
\end{itemize} 

\subsection*{\text{(ε\textlatin{v})}}

\begin{itemize}
    \item \text{Αφού } 20\% (ε) $>$ 1\% (α) $\implies$ \text{η δειγματοληψία (α) συμφέρει περισσότερο και πάλι ...}
\end{itemize} 


\subsection*{\text{(στ\textlatin{i})}}

\begin{itemize}
    \item \text{Πλήθος διακοπών/δευτ.:} $125 \textlatin{(MBytes/s)} / 40 (\textlatin{Bytes}/\text{πακέτο}) = 125 \times 10^6 / 40 = 3.125.000 \times 1\text{ διακοπή/δ}.$
\end{itemize}

\subsection*{\text{(στ\textlatin{ii})}}

\begin{itemize}
    \item \text{Συνολικό κόστος σε κύκλους: } $3.125.000 \times 2000\text{ (κύκλοι/διακοπή)} = 6.250.000.000$ κύκλοι.
\end{itemize}

\subsection*{\text{(στ\textlatin{iii})}}

\begin{itemize}
    \item \text{Επεξεργαστής: } ${6.250.000.000} / {1.000.000.000} \times 100 = 625\% \Rightarrow \text{δεν αντέχει: απαιτήσεις} > \text{δυναμικότητά } ...$
\end{itemize}




\newpage

\text{   }
\begin{center}
    \vspace{-1.5cm}
    {\underline{\Large \textbf{Άσκηση 13.7}}}
\end{center}
\vspace{0.345cm}

\subsection*{\text{Δεδομένα (όλα):}}

\begin{itemize}
    \item Συχνότητα \textlatin{CPU}: $1$~\textlatin{GHz} $= 1$ κύκλος/\textlatin{ns} $= 1.000.000.000$ κύκλοι/δευτερόλεπτο.
    \item Σύνδεσμος \textlatin{PCI-Express x1}: $2.5$~\textlatin{Gbits/s} $= 312$~\textlatin{MBytes/s} $\approx 1$~\textlatin{Byte} κάθε $3.2$ κύκλους ρολογιού.
    \item Μέγεθος πακέτου: $12$~\textlatin{Bytes} \textlatin{(header)} $+ 4$~\textlatin{Bytes} \textlatin{(data)} $= 16$~\textlatin{Bytes}.
    \item Χρόνος μεταφοράς πακέτου: $16 \times 3.2 = 50$ κύκλοι ρολογιού.
    \item Κόστος βρόχου αντιγραφής (εκτός μεταφοράς): $30$ κύκλοι ρολογιού.
    \item Συνολικός χρόνος ανά λέξη (4 \textlatin{Bytes}): $50 + 30 = 80$ κύκλοι ($80$~\textlatin{ns}).
    \item Απαιτήσεις συσκευών:
    \begin{itemize}
        \item $2$ δίσκοι: $10$~\textlatin{MBytes/s} ο καθένας $\Rightarrow 20$~\textlatin{MBytes/s} σύνολο.
        \item $1$ δίκτυο \textlatin{Fast Ethernet}: $100$~\textlatin{Mbits/s} $= 12.5$~\textlatin{MBytes/s}.
    \end{itemize}
    \item Συνολικός ρυθμός δεδομένων: $20 + 12.5 = 32.5$~\textlatin{MBytes/s}.
\end{itemize}

\vspace{0.234567cm}

\subsection*{\text{(α\textlatin{i}) }}

\begin{itemize}
    \item Αντιγραφή: $4$~\textlatin{Bytes} κάθε $80$~\textlatin{ns}.
    \item Κύκλοι των $80$~\textlatin{ns} σε $1$ δευτερόλεπτο: ${10^9} / {80} = 12.5 \times 10^6$.
    \item Ρυθμός σε \textlatin{MBytes/s}: $4 \times 12.5 \times 10^6 = 50 \times 10^6 = 50$~\textlatin{MBytes/s}.
    \item Ρυθμός σε \textlatin{Mbits/s}: $50 \times 8 = 400$~\textlatin{Mbits/s}.
\end{itemize}

\subsection*{\text{(α\textlatin{ii}) }}

\begin{itemize}
    \item Συνολικός ρυθμός δεδομένων: $32.5$~\textlatin{MBytes/s}.
    \item Ρυθμός αντιγραφής επεξεργαστή: $50$~\textlatin{MBytes/s}.
    \item Ποσοστό χρόνου: ${32.5} / {50} \times 100 = 65\%$.
\end{itemize}



\subsection*{\text{(β\textlatin{i}) }}

\begin{itemize}
    \item Μέγεθος πακέτου: $12$~\textlatin{Bytes (header)} $+ 4$~\textlatin{Bytes (data)} $= 16$~\textlatin{Bytes}.
    \item Χρόνος ανά πακέτο: $16 \times 3.2 = 51.2$~\textlatin{ns}.
    \item Πακέτα ανά δευτερόλεπτο: ${10^9} / {51.2} \approx 19.531.250$~\text{πακέτα/s}.
    \item Ρυθμός δεδομένων: $4 \times 19.531.250 = 78.125.000$~\textlatin{Bytes/s} $= 78.125$~\textlatin{MBytes/s}.
    \item Ρυθμός σε \textlatin{Mbits/s}: $78.125 \times 8 = 625$~\textlatin{Mbits/s}.
\end{itemize}

\subsection*{\text{(β\textlatin{ii}) }}

\begin{itemize}
    \item Μέγεθος πακέτου: $12$~\textlatin{Bytes (header)} $+ 64$~\textlatin{Bytes (data)} $= 76$~\textlatin{Bytes}.
    \item Χρόνος ανά πακέτο: $76 \times 3.2 = 243.2$~\textlatin{ns}.
    \item Πακέτα ανά δευτερόλεπτο: ${10^9} / {243.2} \approx 4.111.842$~\text{πακέτα/s}.
    \item Ρυθμός δεδομένων: $64 \times 4.111.842 \approx 263.157.888$~\textlatin{Bytes/s} $\approx 263.158$~\textlatin{MBytes/s}.
    \item Ρυθμός σε \textlatin{Mbits/s}: $263.158 \times 8 \approx 2.105.264$~\textlatin{Mbits/s}.
\end{itemize}

 
\subsection*{\text{(γ\textlatin{i})}}

\begin{itemize} 
    \item Ποσοστό χρόνου απασχόλησης συνδέσμου μνήμης: ${32.5} \text{(α\textlatin{ii})} / {263.158} \text{(β\textlatin{ii})} \times 100 (\%) \approx 12.35\%$.
\end{itemize}

\subsection*{\text{(γ\textlatin{ii})}}

\begin{itemize}
    \item (α) $>$ (γ) $\longrightarrow$ (α\textlatin{ii}) $>$ (γ\textlatin{i}) $\longrightarrow$ 65\% $>$ 12.5\% \ ... 
\end{itemize}

\subsection*{\text{(γ\textlatin{iii})}}

\begin{itemize}
    \item Μεταφορές μεγαλύτερων πακέτων $\implies$ αυξάνεται η παροχή δεδομένων ανά πακέτο \ ...
    \item Εκτελούνται ταυτόχρονα/παράλληλα επιπλέον ανεξάρτητες εντολές $\implies$ μείωση απασχόλησης \ ...
\end{itemize}
  



% ---------------------------------------------------------------------------------------

\newpage

\begin{textblock}{10}(14.401, 15.501)
    \boxed{\textbf{\LaTeX{}}}
\end{textblock}

\begin{center}
    \vspace*{\fill}
    \textbf{\LARGE{ΤΕΛΟΣ}} \\[0.001cm]
    \textbf{\hspace{0.1cm}(ευχαριστώ)}
    \\[0.1cm]
    \vspace*{\fill}
\end{center}

% ---------------------------------------------------------------------------------------

\end{document}

% -------------------------------------------------------------- %
%                           The End :)                           %
% -------------------------------------------------------------- %

